\chapter{Fase 8: la campaña ``El sur tiene espacio''}
\label{cap:d-f8}
\textbf{Fecha}: 2 de octubre de 2026. \textbf{Nota}: \codigo{21-campana.md}.

\textbf{Qué se buscaba}: a quién le habla la campaña, con qué palabras, en qué canal, cuándo y cómo se mide. Cubre los
elementos obligatorios del Problema Prototípico (buyer persona, branding, medios, piezas, implementación e indicadores) y el
incidente de mercadotecnia: promover sin provocar un nuevo colapso.

\begin{opciones}[8.1 · Cuántas personas · 2-oct-2026]
\begin{center}
\small
\begin{tabularx}{\linewidth}{>{\raggedright\arraybackslash}p{0.3\linewidth}Y}
\encabezado \h{Opción} & \h{Por qué sí o por qué no} \\
\textbf{Dos: ``La que vuelve al sur'' y ``La que baja del norte'' (elegida)} & Cada una con su mensaje y su canal \\
\filagris Solo la nacional & Deja fuera al extranjero que ya está a 334 km, en Cancún \\
Solo la extranjera & La Bahía es 95\,\% nacional \\
\bottomrule
\end{tabularx}
\normalsize
\end{center}
\textbf{Evidencia}: Oxtankah 94.8\,\% nacionales y la Ruta 63.3\,\% (INAH 2025); de los 9,408,423 extranjeros en Cancún, 56.3\,\%
de Estados Unidos y 17.1\,\% de Canadá; solo 250 extranjeros llegaron en avión a Chetumal. Cada rasgo de cada persona lleva
su etiqueta: \emph{dato}, \emph{derivado}, \emph{supuesto} o \emph{hueco} (estado de origen, edad e ingreso del nacional).
\end{opciones}

\begin{opciones}[8.2 · El nombre de la campaña · 2-oct-2026]
\textbf{Elegido}: \textbf{``El sur tiene espacio''} (``The south has room''): dice el dato y la promesa a la vez.

\textbf{Descartados}: ``Sur sin prisa'' (no dice dónde) y ``Del Caribe al sur'' (solo le sirve a la extranjera).
\end{opciones}

\begin{opciones}[8.3 · El tono · 2-oct-2026]
\textbf{Opciones presentadas}: cercano y tranquilo; orgullo cultural maya y mexicano; aventura.

\textbf{Lo que eligió Brandon}: \emph{``combina las 3, porque es un orgullo ser méxicano''}.
\end{opciones}

\section*{De dónde salió el mensaje: las reseñas}
De 85,987 reseñas de Quintana Roo (Tulum, Isla Mujeres y Bacalar), el riesgo de que una reseña sea mala sube con el
\textbf{ruido} (2.76 veces), la \textbf{suciedad} (1.89), el \textbf{precio} (1.66) y las \textbf{multitudes} (1.52), y baja con
la \textbf{cultura} (0.48) y la \textbf{calma} (0.41). Lo que hunde una reseña es lo que sobra en los destinos llenos; lo que la
protege es lo que el sur ofrece. Esa es la promesa de la marca.

\begin{opciones}[8.4 · Decisiones técnicas de la minería de texto]
\begin{itemize}
\item \textbf{Temas con un léxico a la vista, no LDA}: así cada tema se define antes y se defiende palabra por palabra.
\item \textbf{Riesgo relativo} como medida (la de la epidemiología): fácil de leer (``2.76 veces más seguido'').
\item \textbf{Reglas de asociación (Apriori)} para combinaciones: ruido + servicio $\Rightarrow$ reseña mala, lift 2.82.
\item \textbf{Log-odds con prior de Dirichlet} para las palabras que distinguen a las reseñas de 5 estrellas
  (\emph{excelente, increíble, hermoso, deliciosa, amable}): evita que las palabras raras salgan por azar.
\item \textbf{Límite declarado}: ninguna reseña es de los cinco lugares.
\end{itemize}
\end{opciones}

\begin{opciones}[8.5 · Reglas de los anuncios]
\begin{itemize}
\item Cada anuncio lleva \textbf{el dato que lo respalda} (por ejemplo, ``Pirámides mayas con espacio'': Tulum 1,031,443
  visitantes, la Ruta 66,628).
\item \textbf{Sin precios ni horas de viaje}, porque no hay dato oficial.
\item El largo se revisa contra los límites de cada plataforma (Google: título $\le30$ y descripción $\le90$; Meta: título
  $\le40$ y texto $\le125$).
\item ``La que baja del norte'' se activa solo cuando Cancún está en temporada alta y el sur tiene espacio, y la Torre la
  enciende por semana.
\end{itemize}
\end{opciones}

\textbf{Resultado}: 2 personas, la marca completa (nombre, lema, personalidad, propuesta de valor, posicionamiento y elementos
visuales), 5 anuncios, el plan de medios (Facebook 70\,\%, Google 30\,\%, WhatsApp sin costo, la página como destino), el
calendario de 9 meses y 10 indicadores (alcance, interacción, conversión, costo por visitante $\le\$196$, afluencia +957,
presupuesto, económico [hueco], social, ambiental y capacidad $\le40\,\%$).

\begin{error}
\begin{itemize}
\item \textbf{Un dominio inventado}: las maquetas de los anuncios mostraban una dirección web que no existe. Se cambió por la
  dirección real de la página, y una prueba lo vigila.
\item \textbf{El calendario} decía que había anuncio para el norte en un mes sin anuncios. Se corrigió.
\end{itemize}
\end{error}
